\section{Iterationen}
\subsection{Slicing}
\begin{itemize}
   \item \textbf{Syntax}: \texttt{part = x[start:end:step]}.
   \item Start ist inklusiv (Standard 0).
   \item End ist exklusiv (Standard End-Index).
   \item Step ist optional (Standard 1). Negativer Step bedeutet umgekehrte Reihenfolge.
\end{itemize}

\begin{lstlisting}[language=Python]
data = [0, 1, 2, 3, 4, 5.0, [1,2,3]]
part = data[::-1] # Umgekehrte Reihenfolge
\end{lstlisting}

\subsection{For-Schleifen}
\begin{itemize}
    \raggedright
    \item Ruft intern \texttt{iter()} auf, um einen Iterator zurückzugeben.
    \item Nutzt die \texttt{next()}-Methode fürs nächste Element, bis die \texttt{StopIteration}-Exception auftritt.
\end{itemize}

\begin{lstlisting}[language=Python]
for value in iterable:
    expression1
\end{lstlisting}

\subsubsection{Schleifensteuerung}
\begin{itemize}
    \item \textbf{break}: Beendet eine Schleife vor ihrer vollständigen Abarbeitung.
    \item \textbf{continue}: Überspringt den restlichen Code der Schleife in der aktuellen Iteration.
    \item \textbf{else-Keyword}: Wird ausgeführt, wenn die Schleife normal ohne ein \texttt{break} endet.
\end{itemize}

\subsubsection{Dictionaries}
\begin{lstlisting}[language=Python]
for key in data:                # Iteriert Schluessel
for value in data.values():     # Iteriert Werte
for key, value in data.items(): # Iteriert Schluessel und Werte
\end{lstlisting}

\subsubsection{Enumerate}
\begin{itemize}
    \item Wird benutzt, wenn beim Iterieren über Elemente ein Zähler benötigt wird.
\end{itemize}

\begin{lstlisting}[language=Python]
for i, fruit in enumerate(fruits, start=10):
    print(i, fruit)
\end{lstlisting}

\subsubsection{Zip}
\begin{itemize}
    \item Iteriert über mehrere Iterables gleichzeitig.
    \item \textbf{Objekte verschiedener Längen}: \texttt{zip()} stoppt, sobald das kürzeste Iterable aufgebraucht ist. Restliche Elemente werden ignoriert.
\end{itemize}

\lstinputlisting[language=Python]{code/zip.py}

\subsection{Comprehensions}
\begin{itemize}
    \item \textbf{Basis-Syntax}: \texttt{[expression for item in iterable]}.
    \item \textbf{Verschachtelte for-Schleifen}: Von links nach rechts ausgewertet.
    \item \textbf{Funktionsaufrufe}: Built-in Funktionen wie \texttt{any()} können benutzt werden.
    \item \textbf{\color{red}Bedingte Syntax}: Die Position der \texttt{if}-Anweisung ändert das Verhalten komplett. 
    \begin{itemize}
        \item \textbf{Filtern (Nur \texttt{if})}: Muss ans \textbf{\color{orange}ENDE} gesetzt werden. Filtert Elemente aus dem Iterable heraus.
        \item \textbf{Modifizieren (\texttt{if/else})}: Muss an den \textbf{\color{orange}ANFANG} (vor \texttt{for}). Ändert den Elementwert bedingungsbasiert, filtert aber nicht.
    \end{itemize}
\end{itemize}

\lstinputlisting[language=Python]{code/comprehensions.py}